\appendix

\section{Proofs}

\subsection{Discrete-Metric Frequency Scoring}
\label{app:exact-match-proof}

\begin{proof}[Proof of Proposition 1]
For any report \(r\),
\[
\E_p[a+b\1\{r=\omega\}]
=
a+b\Pr_p(\omega=r),
\]
so expected-score maximization is equivalent to maximizing the multinomial probability mass.

It remains to characterize the mode region.
Compare \(r\) with \(r+e_i-e_j\), which is feasible when \(r_j>0\).
If \(p_j=0\) and \(r_j>0\), then \(r\) cannot be a mode unless all probability is zero, which is impossible.
For \(p_j>0\),
\[
\frac{\Pr_p(r+e_i-e_j)}{\Pr_p(r)}
=
\frac{r_j}{r_i+1}\frac{p_i}{p_j}.
\]
Thus a necessary condition for \(r\) to be a mode is
\[
r_jp_i\le (r_i+1)p_j
\quad
\forall i,j\text{ with }r_j>0.
\]
Conversely, if these inequalities hold, any feasible \(s\) can be reached from \(r\) by a sequence of one-count transfers from coordinates where the current vector exceeds \(s\) to coordinates where it is below \(s\).
At each step the same ratio is at most one, because source counts weakly decrease and destination counts weakly increase along the path.
Therefore no feasible \(s\) has larger multinomial mass than \(r\).

The coordinate upper bound follows by summing
\[
r_jp_i\le (r_i+1)p_j
\]
over \(j\ne i\):
\[
(n-r_i)p_i\le (r_i+1)(1-p_i),
\]
so \(p_i\le (r_i+1)/(n+1)\).
The lower bound follows by summing
\[
r_ip_j\le (r_j+1)p_i
\]
over \(j\ne i\):
\[
r_i(1-p_i)\le (n-r_i+k-1)p_i,
\]
so \(p_i\ge r_i/(n+k-1)\).

The bounds are sharp.
The upper endpoint is attained by
\[
p_i=\frac{r_i+1}{n+1},
\quad
p_j=\frac{r_j}{n+1}\quad(j\ne i).
\]
The lower endpoint, when \(r_i>0\), is attained by
\[
p_i=\frac{r_i}{n+k-1},
\quad
p_j=\frac{r_j+1}{n+k-1}\quad(j\ne i).
\]
When \(r_i=0\), the lower endpoint \(0\) is attained by setting \(p_i=0\) and assigning probabilities proportional to \(r_j\) on the reported support.
When \(r_i=n\), the upper endpoint \(1\) is attained by \(p_i=1\).
\end{proof}

\subsection{Manhattan-Distance Rule}
\label{app:manhattan-proof}

The Manhattan characterization used in Section 3.2 follows from a standard exchange argument.
Linearity of expectation gives
\[
\E_p\left[\sum_{i=1}^k |r_i-\omega_i|\right]
=
\sum_{i=1}^k \E_p|r_i-\omega_i|.
\]
The \(i\)-th term depends only on the binomial marginal \(W_i\sim\mathrm{Bin}(n,p_i)\).
For
\[
\ell_i(t;p_i)=\E|t-W_i|,
\]
the marginal cost of increasing the reported count from \(t\) to \(t+1\) is
\[
\ell_i(t+1;p_i)-\ell_i(t;p_i)
=
2F_i(t;n,p_i)-1.
\]
Moving one reported count from \(j\) to \(i\) changes expected Manhattan loss by
\[
\{2F_i(r_i;n,p_i)-1\}-\{2F_j(r_j-1;n,p_j)-1\}.
\]
Hence no feasible one-count transfer improves the report if and only if
\[
F_i(r_i;n,p_i)\ge F_j(r_j-1;n,p_j)
\quad
\forall i,j\text{ with }r_i<n,\ r_j>0.
\]
Because the marginal increments are weakly increasing in the reported count, this local no-transfer condition is also sufficient for global optimality on the integer simplex.

The threshold representation follows by taking \(c\) between
\[
\max_{j:r_j>0}F_j(r_j-1;n,p_j)
\quad\text{and}\quad
\min_{i:r_i<n}F_i(r_i;n,p_i).
\]
Inverting the resulting binomial-CDF inequalities coordinate by coordinate gives the threshold-indexed box-simplex description used for Manhattan bounds.

\subsection{Quadratic-Distance Rule}
\label{app:squared-proof}

\begin{proof}[Proof of Proposition 2]
Maximizing \(S_2\) is equivalent to minimizing \(\E_p\|r-\omega\|_2^2\).
For fixed \(r\),
\[
\E_p\|r-\omega\|_2^2
=
\|r-\E_p[\omega]\|_2^2
+
\sum_{i=1}^k\mathrm{Var}_p(\omega_i).
\]
The variance term is independent of \(r\), and \(\E_p[\omega]=np\).
Hence optimal reports are exactly the feasible integer projections of \(np\).

The full inverse projection condition is
\[
\|r-np\|_2^2\le \|s-np\|_2^2
\quad
\forall s\in\OmegaN,
\]
or equivalently
\[
2np\cdot(s-r)\le \|s\|_2^2-\|r\|_2^2.
\]
Taking \(s=r+e_i-e_j\), feasible when \(r_j>0\), gives
\[
n(p_i-p_j)\le r_i-r_j+1.
\]
Conversely, suppose all these one-count inequalities hold.
For any \(s\in\OmegaN\), write \(d=s-r\).
Pair each positive unit of \(d_i\) with a negative unit of \(d_j\).
Summing the one-count inequalities over the paired transfers gives
\[
np\cdot d\le r\cdot d+M,
\]
where \(M=\sum_{i:d_i>0}d_i\) is the number of transferred counts.
Because \(d\) has integer components and sums to zero,
\[
M\le \frac12\sum_i d_i^2.
\]
Therefore
\[
2np\cdot(s-r)
\le
2r\cdot(s-r)+\|s-r\|_2^2
=
\|s\|_2^2-\|r\|_2^2,
\]
which is the full projection condition.

It remains to derive the coordinate bounds.
Let \(m=m(r)=|\{j:r_j>0\}|\).
If \(r_i>0\), summing
\[
n(p_u-p_i)\le r_u-r_i+1
\]
over all \(u\ne i\) gives
\[
1-p_i\le (k-1)p_i+\frac{n-kr_i+k-1}{n},
\]
so
\[
p_i\ge \frac{r_i-1}{n}+\frac{1}{nk}.
\]
For the upper bound, sum
\[
n(p_i-p_v)\le r_i-r_v+1
\]
over positive-report coordinates \(v\ne i\), and use nonnegativity of the zero-report coordinates:
\[
1-p_i\ge (m-1)p_i+\frac{n-mr_i-m+1}{n}.
\]
This gives
\[
p_i\le \frac{r_i+1}{n}-\frac{1}{nm}.
\]

If \(r_i=0\), the lower bound is \(0\).
For the upper bound, sum
\[
n(p_i-p_v)\le 1-r_v
\]
over all \(m\) positive-report coordinates \(v\):
\[
1-p_i\ge mp_i+\frac{n-m}{n},
\]
so
\[
p_i\le \frac{m}{n(m+1)}.
\]
The endpoints are attained by making the relevant summed inequalities bind and assigning remaining probability within the categories not used in the binding constraints.
Thus the displayed bounds are sharp.
\end{proof}

\section{Design Exercise Details}
\label{app:design-details}

The design exercise is generated by the script \texttt{scripts/design\_efficiency.py}.
The final grid uses \(n\in\{5,10,20,50\}\), \(k\in\{2,3,5,10\}\), \(\alpha\in\{0.1,0.3,1,3,10\}\), and \(5{,}000\) draws from the symmetric Dirichlet distribution for each cell.
For each draw and rule, the script computes an optimal report and then computes bounds for the inverse region generated by that report.

Discrete-metric coordinate bounds are closed form and mean bounds are linear programs over the transfer region.
Quadratic-distance coordinate bounds are closed form and mean bounds are linear programs over the one-count transfer polytope.
Manhattan coordinate and mean bounds use the threshold representation in Appendix \ref{app:manhattan-proof}; the script inverts binomial CDF inequalities by generalized inverse functions and searches over a threshold grid with tolerance \(10^{-4}\).
The Manhattan entries are therefore threshold-computed.

The generated files are CSV tables under \texttt{outputs/simulation\_design/}.
These outputs are diagnostics for experimental design.
They are not simulations of subject behavior and do not replace the mathematical characterization of \(R_S(p)\) and \(P_S(r)\).
